\section{Limitations, safety, and conclusion}
\label{sec:limitations}

\paragraph{Approximation boundary.}
The local mixture is not the diffusion posterior.  Jacobian error, finite
differences, PSD/rank truncation, and frozen fantasies add bias; guarantees
hold only for the surrogate.  H2 expansion beams and any root cap are
heuristic.  NFE omits planner arithmetic, so latency and energy remain
essential.

\paragraph{Retrospective sensing.}
Both domains reveal held-out entries only after selection.  This prevents
policy leakage but validates neither scanner timing nor hardware.  LiDAR
budgets count spherical-raster queries, not verified firings, and the
likelihood does not model no-return events.  Nuisance tests remain simulations
until prospective acquisition.

\paragraph{Clinical scope.}
Single-reader, non-exhaustive fastMRI+ regions cannot support diagnostic
claims; PSNR can rise while a lesion is erased.  We therefore require ROI,
detector, hallucination, subgroup, and worst-decile analyses.  Dataset access
agreements still govern redistribution.

\paragraph{Conclusion.}
\method{} combines diffusion modes and local tangents into an
operator-sufficient active-imaging belief.  Its algebra is exact only within
the surrogate; its value depends on real calibration and matched-compute
outcome routing.  The locked MRI/LiDAR study keeps that answer \pending{}
until evidence exists.
